\section{Introduction}
\label{sec:introduction}

\subsection{Why Graph Distance Diverges from Geographic Distance}

Phase~1 Centroidal Voronoi Districts (\cvd{}, B.22~\citep{dellaLibera2026cvd})
assigns each census tract to its nearest seed by BFS hop-count on the
adjacency graph and moves seeds to the graph-distance medoid of their district.
BFS hop-count is a practical proxy for geographic proximity in most inland
states: tracts are roughly uniform in area, and one adjacency hop corresponds
to roughly one tract width of geographic distance.

This approximation breaks down in three state archetypes.

\paragraph{Coastal states.}
Maine, Washington, Florida, and Louisiana have census tracts arranged along
long coastlines with few cross-water adjacencies.
Two tracts on opposite shores of a bay may be 3 miles apart in Euclidean
distance but 80 BFS hops apart along the land corridor.
Phase~1 Voronoi assignment in such states forces coastal tracts into ``wrong''
districts based on graph topology rather than actual geographic nearness.
The result is elongated peninsular districts that cross geographic intuition:
a tract on the eastern shore of Puget Sound is assigned to a Cascade-slope
seed rather than to the geographically closer western-shore seed, because
the land route around the sound is shorter in hop-count than the water
crossing.

\paragraph{Panhandle and elongated states.}
Florida's panhandle extends 200 miles westward from the main peninsula.
In BFS space, panhandle tracts in Escambia County (Pensacola) are far from
panhandle tracts in Jefferson County (Tallahassee) only because the land
corridor runs through many small urban tracts.
Geographic distance along the east-west axis is accurately captured by
Euclidean coordinates; hop-count distance misrepresents it because urban
tracts contribute only 1 hop regardless of geographic span.
Texas presents the same geometry on a larger scale: tracts in El Paso are
geographically 800 miles from Beaumont but the hop-count path through
sparse West Texas rural tracts may undercount the true distance.

\paragraph{Sparse interior states.}
Alaska is the extreme case: vast tundra census tracts are geographically
enormous but each contributes 1 hop in the adjacency graph.
An Alaskan census tract covering 10{,}000 square miles is treated identically
to an Anchorage urban tract covering 2 square miles in hop-count space.
Phase~1 Voronoi seeds placed by BFS in Alaska produce geographic nonsense:
the geometric centre of a district in BFS space may be at the corner of a
city, far from the physical centre of the district's territory.

\subsection{What Phase~2 Fixes}

Phase~2 \cvdgeo{} replaces the BFS hop-count metric with Euclidean distance
on projected census-tract centroids.
The projection is EPSG:5070 Albers Equal Area Conic (continental US), which
maps WGS84 (longitude, latitude) pairs to (x, y) in meters, preserving area
and minimising shape distortion for the redistricting domain~\citep{snyder1987,epsg5070}.

Three consequences follow from this metric change:

\begin{enumerate}
\item \textbf{True centroid update.}
  In Phase~1, the ``centroid'' of a district must be the graph-distance
  medoid --- an actual tract, found by evaluating BFS sums over a probe set.
  In Phase~2, the population-weighted mean of projected (x, y) coordinates
  is a valid Euclidean centroid; the nearest real tract is used as the
  updated seed.
  This is exact Lloyd's algorithm on the projected point set, without
  approximation.

\item \textbf{Faster convergence.}
  The Euclidean energy surface is smoother than the BFS hop-count surface.
  Empirically, Phase~2 converges in 7--12 iterations on the test states,
  versus 11--15 iterations for Phase~1.

\item \textbf{No BFS evaluations for medoid.}
  Phase~1 requires $2 \times 50$ single-source BFS evaluations per iteration
  to compute the approximate medoid.
  Phase~2 replaces this with $O(m)$ arithmetic for the weighted mean
  and a nearest-neighbour scan.
  Runtime is strictly faster than Phase~1 for the same iteration count.
\end{enumerate}

\subsection{Scope and Main Results}

This paper is structured as a companion to B.22~\citep{dellaLibera2026cvd}.
Section~\ref{sec:comparison-overview} gives a direct Phase~1 vs.\ Phase~2
property comparison.
Section~\ref{sec:algorithm} gives the full Phase~2 algorithm with pseudocode,
including the Albers projection formula, population-weighted centroid computation,
nearest-tract snap, and $\varepsilon$-convergence.
Section~\ref{sec:results} presents empirical Polsby-Popper comparisons on
NC ($k=14$), FL ($k=28$), and WA ($k=10$) --- spanning the spectrum from
regular geometry to severely coastal.
Section~\ref{sec:data-pipeline} describes the centroid companion file format
and the \texttt{redist fetch --type centroids} pipeline step.
Section~\ref{sec:conclusion} gives recommendations for when to use
\texttt{--cvd-metric geographic} vs.\ the Phase~1 default.

\paragraph{Main result.}
Phase~2 \cvdgeo{} is a strict improvement over Phase~1 \cvd{} for coastal and
elongated states: Florida gains $\approx +6\%$ Polsby-Popper; Washington gains
$\approx +4\%$; North Carolina (regular geometry) gains $\approx +2\%$.
The improvement comes at no runtime cost --- Phase~2 is faster per iteration
than Phase~1 because the medoid computation is replaced by a single
population-weighted mean.

\subsection{Relation to Prior Work}

The continuous CVT framework of Du, Faber, and Gunzburger~\citeyear{du1999}
applies Lloyd's algorithm in Euclidean space.
Phase~2 \cvdgeo{} is the direct discretisation of that framework to
census-tract redistricting: the Euclidean metric is realised via EPSG:5070
projection, and the continuous centroid is approximated by the nearest real tract.
Phase~1 \cvd{} is the graph-distance adaptation; Phase~2 \cvdgeo{} is the
Euclidean adaptation.

Snyder~\citeyear{snyder1987} provides the definitive reference for the Albers
Equal-Area Conic projection used here.
The spherical-Albers formula (our implementation) introduces at most 0.3\%
geometric error versus the full GRS80 ellipsoidal formula; at census-tract
resolution this is $\le 3$\,m per tract, negligible for redistricting purposes.
